% ============================================================
% SECTION 6 — Universal Justice & the Promised al-Mahdi (Chapter 8)
% ============================================================
\section{Justice \& Mahdism}

\begin{frame}{What Is Justice?}
\begin{itemize}[itemsep=5pt]
  \item Prophetic missions serve two aims: right relation to God (tawhid), right relation to each other (justice)
  \item Justice = giving each person their due based on nature, qualifications, and actions -- \emph{not} mechanically identical treatment
  \item Universal justice = the eventual end of tyranny, discrimination, war, and exploitation -- realized in \emph{this} world, not deferred to the next
\end{itemize}
\end{frame}

\begin{frame}{Three Rivals, and 'Ali as the Counter-Example}
\begin{center}
\begin{tikzpicture}[every node/.style={font=\scriptsize,align=center},
  rival/.style={rectangle,rounded corners,draw=softgray,fill=white,text width=3.4cm,minimum height=1.6cm},
  islam/.style={rectangle,rounded corners,draw=deepteal,thick,fill=lightbg,text width=4.2cm,minimum height=1.6cm}]
\node[rival] (n) {\textbf{Nietzsche/Machiavelli}\\ power is fundamental; justice = invention of the weak};
\node[rival, right=0.5cm of n] (r) {\textbf{Russell}\\ justice = enlightened self-interest};
\node[rival, right=0.5cm of r] (m) {\textbf{Marxist}\\ justice emerges automatically from historical/economic development};
\node[islam, below=1cm of r] (i) {\textbf{Islamic view}\\ people can be \emph{educated} to want justice intrinsically -- 'Ali as living proof};
\draw[-{Latex},thick,gold] (n) -- (i);
\draw[-{Latex},thick,gold] (r) -- (i);
\draw[-{Latex},thick,gold] (m) -- (i);
\end{tikzpicture}
\end{center}
\bigquote{Are you asking me to gain victory and success in politics at the price of oppression?}{'Ali -- the same refusal from Chapter 1, now doing philosophical work}
\delivernote{Russell's self-interest theory is the one worth dwelling on: it predicts commitment to justice should weaken once someone is powerful enough that retaliation is no longer a threat. Ask the room whether that's an acceptable outcome.}
\end{frame}

\begin{frame}{Mahdism: A History, Not an Invention}
\begin{center}
\begin{tikzpicture}[every node/.style={font=\tiny,align=center},
  node distance=0.35cm,
  flowbox/.style={rectangle,rounded corners,draw=deepteal,fill=lightbg,text width=2.5cm,minimum height=1.3cm}]
\node[flowbox] (a) {'Ali to Kumayl:\\earth never without a proof, hidden or manifest};
\node[flowbox, right=of a] (b) {Mukhtar presents Ibn Hanafiyyah as Mahdi};
\node[flowbox, right=of b] (c) {Followers: hidden at Mt.\ Radwa, not dead};
\node[flowbox, right=of c] (d) {Nafs al-Zakiyyah floated at Abwa' -- al-Sadiq rejects the claim, not the cause};
\node[flowbox, right=of d] (e) {Mansur names his son ``Mahdi'' -- privately doubts any claimant};
\foreach \x/\y in {a/b,b/c,c/d,d/e}{\draw[-{Latex},thick,gold] (\x) -- (\y);}
\end{tikzpicture}
\end{center}
\begin{itemize}[itemsep=3pt]
  \item Belief predates Imam al-Mahdi's own birth by well over a century
  \item Documented in Sunni sources too, not only Shi'ah ones -- even an Abbasid caliph was named ``al-Mahdi''
  \item Al-Sadiq publicly demonstrated Isma'il's death specifically to prevent a false Mahdi claim forming around him
\end{itemize}
\end{frame}

\begin{frame}{The Central Caution}
\begin{center}
\Large\bfseries\color{deepteal}
The existence of a claim does not prove the claim.
\end{center}
\vspace{0.3cm}
\examplenote{Muhammad ibn 'Ijlan}{A Medinan scholar pledged allegiance to 'Abd Allah al-Mahd's son believing him the promised Mahdi -- other scholars had to intervene on his behalf. Real, tangible stakes for getting this wrong.}
\end{frame}

\begin{frame}{Awaiting Is Not Passivity}
\bigquote{Awaiting does not invalidate obligatory duties.}{Mutahhari}
\begin{itemize}[itemsep=5pt]
  \item The dangerous inference: if reappearance follows corruption, maybe corruption should be encouraged
  \item The ``abscess'' analogy explains, then explicitly rejects, that revolutionary model
  \item Prayer, fasting, zakat, Hajj, enjoining good, forbidding evil -- \emph{all} remain fully in force
  \item Correct awaiting is active: fulfilling duty, training society, pursuing reform now
\end{itemize}
\end{frame}

\qaframe{Justice \& Mahdism}{%
  \qa{Which of the book's own historical episodes would each rival theory of justice struggle hardest to explain?}
     {Power theory struggles with 'Ali's repeated refusal of advantage (Ch.\ 1); Russell's self-interest theory struggles with Musa al-Kazim's refusal of a costless confession (Ch.\ 5); the Marxist theory struggles to explain any individual moral choice at all, since it locates the engine of justice in impersonal forces rather than personal virtue.}
  \qa{Why is the ``abscess'' rejection arguably the single most practically important claim in the entire book?}
     {It's the point where doctrine becomes dangerously actionable: if corruption is a precondition for reappearance, hastening it could seem to hasten deliverance. Mutahhari's explicit rejection closes off the most consequential possible misreading of everything the book has built.}
  \qa{Why does Chapter 8 function as the book's synthesis point rather than the Introduction?}
     {The Introduction supplies the interpretive tool (essence/expression) but not yet the material to apply it to. Chapter 8 arrives after every Imam's specific circumstance has been presented, so it can draw the threads together: justice as the shared essence, Mahdism as its promised full expression.}
}
